\chapter{矩阵分析}
\label{chap:matrix-analysis}
线性代数告诉我们怎样表示线性映射、怎样解线性方程组，但机器学习还会追问一些更细的
问题。两个特征几乎重复时，为什么最小二乘系数会对微小噪声极其敏感？压缩一个大矩阵时，
保留哪些方向才能使误差最小？计算机返回数值解以后，怎样判断误差来自数据还是求解算法？
这些问题不能只用“矩阵是否可逆”回答，还要知道矩阵沿不同方向放大多少、哪些方向几乎
被压扁，以及这种结构在扰动下是否稳定。
本章继续使用第~\ref{chap:linear-algebra}章的近共线样本。把输出空间换到
\(\bm{c}_1/\sqrt{3}\)与\(\bm{q}/\sqrt{2}\)组成的标准正交坐标，并将整个输出统一
除以\(\sqrt{3}\)，原设计矩阵化为下面的二维表示，其中
\(\widetilde{\varepsilon}=\sqrt{2/3}\,\varepsilon\)：
\begin{equation}
  \bm{X}_{\widetilde{\varepsilon}}=
  \begin{bmatrix}1&1\\0&\widetilde{\varepsilon}\end{bmatrix},
  \qquad 0<\widetilde{\varepsilon}\ll 1.
  \label{eq:matrix-analysis-nearly-collinear-design}
\end{equation}
它的两列是$\bm{x}_1=(1,0)^{\top}$与
$\bm{x}_2=(1,\widetilde{\varepsilon})^{\top}$，在$\widetilde{\varepsilon}$很小时几乎同向。对响应
$\bm{y}=(2,0)^{\top}$，方程$\bm{X}_{\widetilde{\varepsilon}}\bm{\beta}=\bm{y}$的解是
$\bm{\beta}=(2,0)^{\top}$；若把响应改为
$\bm{y}+\Delta\bm{y}=(2,\eta)^{\top}$，解就变成
\begin{equation}
  \bm{\beta}+\Delta\bm{\beta}=
  \begin{bmatrix}2-\eta/\widetilde{\varepsilon}\\ \eta/\widetilde{\varepsilon}\end{bmatrix}.
  \label{eq:matrix-analysis-collinear-coefficient-perturbation}
\end{equation}
响应只改变$|\eta|$，两个系数却分别改变$|\eta|/\widetilde{\varepsilon}$，而预测向量只从
$(2,0)^{\top}$变为$(2,\eta)^{\top}$。参数敏感与预测敏感不是同一件事；“几乎
重复的列”背后存在一个被矩阵强烈压缩的方向。

我们先用特征值和奇异值找出这些方向，再讨论低秩近似怎样舍弃不可靠方向。矩阵范数将
方向性结论汇总成尺度，矩阵函数与结构化恒等式帮助操纵大型线性变换；条件数和扰动理论
区分问题敏感性与算法误差，稳定求解方法则把这些认识落实为计算策略。

\section{谱揭示矩阵沿哪些方向作用}
\label{sec:matrix-analysis-spectrum}

\subsection{特征值描述不改变方向的缩放}

给定方阵$\bm{A}\in\mathbb{R}^{d\times d}$，若存在非零向量
$\bm{v}\in\mathbb{C}^d$和标量$\lambda\in\mathbb{C}$使
\begin{equation}
  \bm{A}\bm{v}=\lambda\bm{v}.
  \label{eq:matrix-analysis-eigenpair}
\end{equation}
则称$\lambda$是$\bm{A}$的\term{特征值}（eigenvalue），$\bm{v}$是对应的
\term{特征向量}（eigenvector）。这意味着$\bm{A}$保持$\bm{v}$张成的一维复子空间
不变，并把长度缩放为原来的$|\lambda|$倍。对实特征值，$\lambda<0$还会使方向反向；
对非实特征值，作用中还包含复相位的变化。实矩阵也可能有复特征值，因此定义允许在
$\mathbb{C}$中寻找特征对。

式~\eqref{eq:matrix-analysis-eigenpair}等价于
$(\bm{A}-\lambda\bm{I})\bm{v}=\bm{0}$。非零解存在当且仅当
\begin{equation}
  p_{\bm{A}}(\lambda)=\det(\lambda\bm{I}-\bm{A})=0.
  \label{eq:matrix-analysis-characteristic-polynomial}
\end{equation}
其中$p_{\bm{A}}$称为特征多项式。一个特征值作为多项式根出现的次数是它的代数重数，
对应特征空间
$\{\bm{v}:\bm{A}\bm{v}=\lambda\bm{v}\}$的维数是几何重数。几何重数不超过代数
重数；两者不等时，矩阵可能没有足够多的线性无关特征向量。

若$\bm{A}$有$d$个线性无关特征向量，把它们按列组成$\bm{V}$，把对应特征值放入
对角矩阵$\bm{\Lambda}$，便有
\[
  \bm{A}\bm{V}=\bm{V}\bm{\Lambda},
  \qquad \bm{A}=\bm{V}\bm{\Lambda}\bm{V}^{-1}.
\]
这称为特征分解或对角化。对角形式使幂运算变得直接：
$\bm{A}^k=\bm{V}\bm{\Lambda}^k\bm{V}^{-1}$。但$\bm{V}$若接近奇异，坐标变换本身
会放大误差；所以非对称矩阵即使可以对角化，也不一定适合直接用特征向量进行数值计算。

即使矩阵不能对角化，它的高次幂也不是彼此独立的新对象。
\term{Cayley--Hamilton定理}（Cayley--Hamilton theorem）断言，每个方阵都满足
自己的特征多项式。若
\[
  p_{\bm{A}}(t)=t^d+c_{d-1}t^{d-1}+\cdots+c_1t+c_0,
\]
则
\[
  p_{\bm{A}}(\bm{A})
  =\bm{A}^d+c_{d-1}\bm{A}^{d-1}+\cdots+c_1\bm{A}+c_0\bm{I}
  =\bm{0}.
\]
因此$\bm{A}^d$可以写成
$\bm{I},\bm{A},\ldots,\bm{A}^{d-1}$的线性组合；反复代入后，任意更高次幂也可以。
这个结论不要求$\bm{A}$可对角化。第~\ref{chap:control-theory}章分析线性系统的可达空间时，
将用它说明超过状态维数的矩阵幂不会产生新的可达方向。

\subsection{实对称矩阵具有正交谱分解}

机器学习中常见的协方差矩阵、Gram矩阵和Hessian矩阵都是实对称矩阵。对称性带来比一般
方阵强得多的结论\scite{math-horn2013matrix}

\begin{theorem}[实对称矩阵的谱定理]
\label{thm:matrix-analysis-symmetric-spectral-theorem}
若$\bm{A}\in\mathbb{R}^{d\times d}$且$\bm{A}^{\top}=\bm{A}$，则它的全部
特征值都是实数，并存在正交矩阵$\bm{Q}$与实对角矩阵$\bm{\Lambda}$使
\begin{equation}
  \bm{A}=\bm{Q}\bm{\Lambda}\bm{Q}^{\top}.
  \label{eq:matrix-analysis-symmetric-eigendecomposition}
\end{equation}
\end{theorem}

\begin{proof}
先把$\bm{A}$看作复矩阵。若
$\bm{A}\bm{v}=\lambda\bm{v}$且$\bm{v}\neq\bm{0}$，则
$\bm{v}^{*}\bm{A}\bm{v}=\lambda\bm{v}^{*}\bm{v}$。左侧因
$\bm{A}^{*}=\bm{A}$而为实数，故$\lambda$为实数；把$\bm{v}$拆成实部与虚部，
其中至少一个给出非零实特征向量。不同特征值对应的向量满足
\[
  \lambda_j\bm{v}_i^{\top}\bm{v}_j
  =\bm{v}_i^{\top}\bm{A}\bm{v}_j
  =(\bm{A}\bm{v}_i)^{\top}\bm{v}_j
  =\lambda_i\bm{v}_i^{\top}\bm{v}_j,
\]
所以二者正交。

对维数归纳。取一个实单位特征向量$\bm{q}_1$。若
$\bm{q}_1^{\top}\bm{x}=0$，则
$\bm{q}_1^{\top}\bm{A}\bm{x}
=(\bm{A}\bm{q}_1)^{\top}\bm{x}=0$，故其正交补在$\bm{A}$下不变，且限制后的
算子仍对称。归纳假设在该$(d-1)$维子空间给出标准正交特征基；与$\bm{q}_1$合并
即可组成$\bm{Q}$，并得到式~\eqref{eq:matrix-analysis-symmetric-eigendecomposition}。
\end{proof}

谱定理中的$\bm{Q}^{-1}=\bm{Q}^{\top}$，所以变换到特征坐标不会改变欧氏长度。
这正是对称特征分解数值行为较好的根本原因。若
$\bm{x}=\sum_{i=1}^{d}c_i\bm{q}_i$，则
\begin{equation}
  \bm{x}^{\top}\bm{A}\bm{x}=\sum_{i=1}^{d}\lambda_i c_i^2.
  \label{eq:matrix-analysis-quadratic-spectral-form}
\end{equation}
因此，对称矩阵为正定，当且仅当全部特征值为正；为半正定，当且仅当全部特征值非负。
复数域中的对应对象是Hermitian矩阵
$\bm{A}^{*}=\bm{A}$，其中$*$表示共轭转置；把上面的转置换成共轭转置后，谱定理与
正定性的谱判据仍成立。一般复矩阵只满足普通转置对称并不足以保证这些结论。

回到式~\eqref{eq:matrix-analysis-nearly-collinear-design}。它的Gram矩阵为
\[
  \bm{G}_{\widetilde{\varepsilon}}=\bm{X}_{\widetilde{\varepsilon}}^{\top}\bm{X}_{\widetilde{\varepsilon}}
  =\begin{bmatrix}1&1\\1&1+\widetilde{\varepsilon}^2\end{bmatrix}.
\]
由特征方程可得
\begin{equation}
  \lambda_{\pm}
  =\frac{2+\widetilde{\varepsilon}^2\mathbin{\pm}\sqrt{4+\widetilde{\varepsilon}^4}}{2}.
  \label{eq:matrix-analysis-nearly-collinear-gram-eigenvalues}
\end{equation}
当$\widetilde{\varepsilon}\to0$时，
$\lambda_+\to2$而$\lambda_-\sim\widetilde{\varepsilon}^2/2$。大特征值对应两列共同变化的
方向，小特征值对应近似相消的方向。后者正是系数可以大幅改变、预测却只改变一点的原因。

\subsection{Rayleigh商把方向优化变成特征值问题}

对实对称矩阵$\bm{A}$和非零向量$\bm{x}$，定义
\term{Rayleigh商}（Rayleigh quotient）
\begin{equation}
  R_{\bm{A}}(\bm{x})
  =\frac{\bm{x}^{\top}\bm{A}\bm{x}}{\bm{x}^{\top}\bm{x}}.
  \label{eq:matrix-analysis-rayleigh-quotient}
\end{equation}
它衡量二次型在方向$\bm{x}$上的强度，并且不受$\bm{x}$长度影响。

\begin{theorem}[Rayleigh--Ritz变分原理]
\label{thm:matrix-analysis-rayleigh-ritz}
设实对称矩阵$\bm{A}$的特征值按
$\lambda_1\geq\cdots\geq\lambda_d$排列，则
\begin{equation}
  \lambda_d=\min_{\bm{x}\neq\bm{0}}R_{\bm{A}}(\bm{x}),
  \qquad
  \lambda_1=\max_{\bm{x}\neq\bm{0}}R_{\bm{A}}(\bm{x}).
  \label{eq:matrix-analysis-rayleigh-ritz}
\end{equation}
等号分别在最小、最大特征值对应的非零特征向量上取得。
\end{theorem}

\begin{proof}
用谱定理写
$\bm{x}=\sum_i c_i\bm{q}_i$。由正交性，
\[
  R_{\bm{A}}(\bm{x})=\frac{\sum_i\lambda_i c_i^2}{\sum_i c_i^2}.
\]
右端是特征值的加权平均，权重
$c_i^2/\sum_jc_j^2$非负且总和为$1$，所以它位于
$[\lambda_d,\lambda_1]$内。取$\bm{x}=\bm{q}_d$或$\bm{q}_1$分别达到两端。
\end{proof}

Weyl不等式需要逐个刻画中间特征值，而不只是谱的两个端点。为此使用
Courant--Fischer极小极大原理的如下版本\scite{math-horn2013matrix}

\begin{theorem}[Courant--Fischer极小极大原理]
\label{thm:matrix-analysis-courant-fischer}
设实对称矩阵$\bm{A}\in\mathbb{R}^{d\times d}$的特征值按
$\lambda_1\geq\cdots\geq\lambda_d$排列。对$i=1,\ldots,d$，
\begin{equation}
  \lambda_i
  =
  \max_{\dim S=i}\ \min_{\substack{\bm{x}\in S\\\bm{x}\neq\bm{0}}}
  R_{\bm{A}}(\bm{x})
  =
  \min_{\dim T=d-i+1}\ \max_{\substack{\bm{x}\in T\\\bm{x}\neq\bm{0}}}
  R_{\bm{A}}(\bm{x}).
  \label{eq:matrix-analysis-courant-fischer}
\end{equation}
\end{theorem}

\begin{proof}
记$\bm{A}$的一组标准正交特征向量为
$\bm{q}_1,\ldots,\bm{q}_d$，并令
$E_i=\operatorname{span}\{\bm{q}_1,\ldots,\bm{q}_i\}$、
$F_i=\operatorname{span}\{\bm{q}_i,\ldots,\bm{q}_d\}$。
任意$i$维子空间$S$都与$F_i$有非零交向量；该向量的Rayleigh商不超过
$\lambda_i$，所以第一个极值不超过$\lambda_i$。取$S=E_i$时，其中每个非零向量的
Rayleigh商都不小于$\lambda_i$，故等号成立。

同理，任意$(d-i+1)$维子空间$T$都与$E_i$有非零交向量；该向量的Rayleigh商不小于
$\lambda_i$，所以第二个极值不小于$\lambda_i$。取$T=F_i$时，其中每个非零向量的
Rayleigh商都不超过$\lambda_i$，再次得到等号。
\end{proof}

这个定理不只是计算特征值的另一种写法。它说明最大特征向量是使二次型最强的单位方向，
最小特征向量则是最弱方向。对
$\bm{G}_{\widetilde{\varepsilon}}=\bm{X}_{\widetilde{\varepsilon}}^{\top}\bm{X}_{\widetilde{\varepsilon}}$，
\[
  R_{\bm{G}_{\widetilde{\varepsilon}}}(\bm{z})=
  \frac{\lVert\bm{X}_{\widetilde{\varepsilon}}\bm{z}\rVert_2^2}
       {\lVert\bm{z}\rVert_2^2}.
\]
因此$\lambda_-$直接量化$\bm{X}_{\widetilde{\varepsilon}}$对最弱系数方向的压缩程度，而不只是
特征多项式的一个根。

只找一个最强方向还不足以决定应保留哪个$k$维子空间。令
$\bm{W}\in\mathbb{R}^{d\times k}$满足
$\bm{W}^{\top}\bm{W}=\bm{I}$，则子空间在二次型下保留的总强度为
\[
  \operatorname{tr}(\bm{W}^{\top}\bm{A}\bm{W}).
\]
把$\bm{W}$的各列展开到$\bm{A}$的正交特征基，并交换两次求和，可得
\[
  \max_{\bm{W}^{\top}\bm{W}=\bm{I}}
  \operatorname{tr}(\bm{W}^{\top}\bm{A}\bm{W})
  =\sum_{i=1}^{k}\lambda_i.
\]
最大值由前$k$个特征向量张成的子空间达到。这是Rayleigh--Ritz原理的子空间形式：
固定维数以后，应保留谱最大的方向。若边界处
$\lambda_k=\lambda_{k+1}$，最优值仍确定，但可以在对应重特征值子空间内更换标准
正交基，甚至有多个最优$k$维子空间；结论不能加强成唯一的特征向量选择。

\subsection{广义特征值比较两个二次型}

有时需要比较的不是
$\bm{x}^{\top}\bm{A}\bm{x}$与欧氏长度，而是两个二次型。设
$\bm{A},\bm{B}\in\mathbb{R}^{d\times d}$，其中
$\bm{A}^{\top}=\bm{A}$且$\bm{B}^{\top}=\bm{B}\succ\bm{0}$。
\term{广义特征值问题}（generalized eigenvalue problem）写作
\begin{equation}
  \bm{A}\bm{v}=\lambda\bm{B}\bm{v}.
  \label{eq:matrix-analysis-generalized-eigenproblem}
\end{equation}
对应的广义Rayleigh商是
\[
  R_{\bm{A},\bm{B}}(\bm{x})
  =\frac{\bm{x}^{\top}\bm{A}\bm{x}}{\bm{x}^{\top}\bm{B}\bm{x}}.
\]
由于$\bm{B}$正定，存在可逆平方根$\bm{B}^{1/2}$。令
$\bm{z}=\bm{B}^{1/2}\bm{x}$，便把广义Rayleigh商改写为
\[
  \frac{\bm{z}^{\top}
    \bm{B}^{-1/2}\bm{A}\bm{B}^{-1/2}\bm{z}}
       {\bm{z}^{\top}\bm{z}}.
\]
同一变量变换还把广义特征方程化为
\[
  \bm{B}^{-1/2}\bm{A}\bm{B}^{-1/2}\bm{z}
  =\lambda\bm{z}.
\]
中间矩阵是实对称矩阵，因此普通Rayleigh--Ritz结论可以直接应用。若广义特征值按
$\lambda_1\geq\cdots\geq\lambda_d$排列，则
\[
  \lambda_d
  =\min_{\bm{x}^{\top}\bm{B}\bm{x}=1}
    \bm{x}^{\top}\bm{A}\bm{x},
  \qquad
  \lambda_1
  =\max_{\bm{x}^{\top}\bm{B}\bm{x}=1}
    \bm{x}^{\top}\bm{A}\bm{x}.
\]
普通特征向量$\bm{z}$通过
$\bm{x}=\bm{B}^{-1/2}\bm{z}$映回广义特征向量。这里不宜显式计算
$\bm{B}^{-1/2}$；实际算法通常用Cholesky分解完成等价变量变换。

$\bm{B}$只半正定时，分母可能为零，商在某些非零方向上没有定义；$\bm{B}$不定时，
分母还会改变符号。此时广义特征值可能包含无穷特征值或复特征值，不能照搬正定情形的
最值解释。条件“$\bm{B}$正定”不是技术装饰，而是把第二个二次型变成合法长度的保证。

\section{奇异值分解连接几何、低秩与最小二乘}
\label{sec:matrix-analysis-svd-low-rank}

\subsection{任意矩阵都有正交的输入与输出方向}

特征分解要求方阵，并且一般矩阵未必可对角化。对任意
$\bm{A}\in\mathbb{R}^{m\times n}$，\term{奇异值分解}
（singular value decomposition, SVD）都给出
\begin{equation}
  \bm{A}=\bm{U}\bm{\Sigma}\bm{V}^{\top},
  \label{eq:matrix-analysis-svd}
\end{equation}
其中$\bm{U}\in\mathbb{R}^{m\times m}$与
$\bm{V}\in\mathbb{R}^{n\times n}$是正交矩阵，
$\bm{\Sigma}\in\mathbb{R}^{m\times n}$只有主对角线上可能非零。其非零对角元按
\[
  \sigma_1\geq\sigma_2\geq\cdots\geq\sigma_r>0
\]
排列，称为奇异值，$r=\operatorname{rank}(\bm{A})$。
上式是完整SVD。只保留前$r$个非零奇异方向得到紧致SVD
$\bm{A}=\bm{U}_r\bm{\Sigma}_r\bm{V}_r^{\top}$，其中三个因子的形状依次为
$m\times r$、$r\times r$与$r\times n$。令$q=\min(m,n)$，经济型或薄SVD则保留
$q$个方向，三个因子的形状依次为$m\times q$、$q\times q$与$q\times n$，其中
仍可能包含零奇异值。写出采用哪一种形式可以避免把人为补齐的基向量误当成数据确定的
方向。

这个分解可以从对称矩阵
$\bm{A}^{\top}\bm{A}$得到。它半正定，因而有标准正交特征向量
$\bm{v}_i$与非负特征值$\sigma_i^2$。对$\sigma_i>0$，定义
$\bm{u}_i=\bm{A}\bm{v}_i/\sigma_i$，则
\[
  \bm{u}_i^{\top}\bm{u}_j=
  \frac{\bm{v}_i^{\top}\bm{A}^{\top}\bm{A}\bm{v}_j}
       {\sigma_i\sigma_j}=\delta_{i,j}.
\]
补全两侧的标准正交基便得到式~\eqref{eq:matrix-analysis-svd}。只保留非零奇异值及其
对应向量，可写成紧致形式
\begin{equation}
  \bm{A}=\sum_{i=1}^{r}\sigma_i\bm{u}_i\bm{v}_i^{\top}.
  \label{eq:matrix-analysis-svd-rank-one-sum}
\end{equation}

SVD给出了明确的几何过程：先用$\bm{V}^{\top}$旋转或反射输入坐标，再沿正交坐标轴
分别按$\sigma_i$缩放，最后用$\bm{U}$旋转或反射到输出空间。零奇异值对应被压到零的
输入方向。由此立即得到
\[
  \operatorname{null}(\bm{A})=
  \operatorname{span}\{\bm{v}_{r+1},\ldots,\bm{v}_n\},
  \qquad
  \operatorname{range}(\bm{A})=
  \operatorname{span}\{\bm{u}_1,\ldots,\bm{u}_r\}.
\]

对贯穿本章的$\bm{X}_{\widetilde{\varepsilon}}$，奇异值满足
$\sigma_{\pm}^2=\lambda_{\pm}$，其中$\lambda_{\pm}$由
式~\eqref{eq:matrix-analysis-nearly-collinear-gram-eigenvalues}给出。因此
\[
  \sigma_{\max}\to\sqrt{2},
  \qquad
  \sigma_{\min}\sim\frac{\widetilde{\varepsilon}}{\sqrt{2}}.
\]
近共线性不是“两个数字看起来接近”这么简单，而是存在一个单位系数方向，被设计矩阵缩小
到约$\widetilde{\varepsilon}/\sqrt{2}$倍。
当若干奇异值相等时，对应左右奇异向量可以同步旋转，单个向量并不唯一；由这些向量
张成的左右奇异子空间才是不依赖基选择的对象。这个区别将在扰动分析中决定应该比较
向量还是子空间。

\subsection{截断SVD给出最佳低秩近似}

若大部分奇异值很小，式~\eqref{eq:matrix-analysis-svd-rank-one-sum}提示我们只保留
前$k$项。定义
\begin{equation}
  \bm{A}_k=\sum_{i=1}^{k}\sigma_i\bm{u}_i\bm{v}_i^{\top},
  \qquad k<r.
  \label{eq:matrix-analysis-truncated-svd}
\end{equation}
它的秩至多为$k$。问题在于，按奇异值截断是否真的比其他秩$k$矩阵更好
\scite{eckart1936}

以下先局部约定矩阵$\bm{M}$的\term{谱范数}（spectral norm）为
\[
  \lVert\bm{M}\rVert_2
  =
  \max_{\bm{x}\neq\bm{0}}
  \frac{\lVert\bm{M}\bm{x}\rVert_2}{\lVert\bm{x}\rVert_2}
  =
  \sigma_1(\bm{M}).
\]
它衡量矩阵对单位输入的最大放大倍数；第
~\ref{sec:matrix-analysis-norms-effective-rank}节将把它与其他矩阵范数统一比较。

\begin{theorem}[Eckart--Young--Mirsky定理]
\label{thm:matrix-analysis-best-low-rank}
设$\bm{A}$的奇异值为
$\sigma_1\geq\cdots\geq\sigma_r>0$。对任意
$0\leq k<r$，
\begin{align}
  \min_{\operatorname{rank}(\bm{B})\leq k}
  \lVert\bm{A}-\bm{B}\rVert_2
  &=\sigma_{k+1},
  \label{eq:matrix-analysis-best-rank-k-spectral}\\
  \min_{\operatorname{rank}(\bm{B})\leq k}
  \lVert\bm{A}-\bm{B}\rVert_{\mathrm{F}}
  &=\left(\sum_{i=k+1}^{r}\sigma_i^2\right)^{1/2}.
  \label{eq:matrix-analysis-best-rank-k-frobenius}
\end{align}
两个最小值都由截断SVD矩阵$\bm{A}_k$达到。
\end{theorem}

\begin{proof}
先证谱范数结论。任取秩不超过$k$的$\bm{B}$。子空间
$S=\operatorname{span}\{\bm{v}_1,\ldots,\bm{v}_{k+1}\}$维数为$k+1$，
而$\operatorname{null}(\bm{B})$的维数至少为$n-k$。维数公式保证
$S\cap\operatorname{null}(\bm{B})$含有非零向量。取其中单位向量
$\bm{x}=\sum_{i=1}^{k+1}c_i\bm{v}_i$，则
\[
  \lVert(\bm{A}-\bm{B})\bm{x}\rVert_2
  =\lVert\bm{A}\bm{x}\rVert_2
  =
  \left(\sum_{i=1}^{k+1}\sigma_i^2c_i^2\right)^{1/2}
  \geq\sigma_{k+1}.
\]
所以任何候选$\bm{B}$的误差都不小于$\sigma_{k+1}$。另一方面，
$\bm{A}-\bm{A}_k=\sum_{i>k}\sigma_i\bm{u}_i\bm{v}_i^{\top}$，
其最大奇异值恰为$\sigma_{k+1}$，下界可以达到。

再证Frobenius范数结论。正交变换不改变Frobenius范数，令
$\bm{C}=\bm{U}^{\top}\bm{B}\bm{V}$，则
\[
  \lVert\bm{A}-\bm{B}\rVert_{\mathrm{F}}^2
  =\lVert\bm{\Sigma}-\bm{C}\rVert_{\mathrm{F}}^2.
\]
把$\bm{C}$的列空间投影记为$\bm{P}$。由于
$\operatorname{rank}(\bm{P})\leq k$且$\bm{P}\bm{C}=\bm{C}$，
\[
  \lVert\bm{\Sigma}-\bm{C}\rVert_{\mathrm{F}}^2
  =\lVert(\bm{I}-\bm{P})\bm{\Sigma}\rVert_{\mathrm{F}}^2+
  \lVert\bm{P}\bm{\Sigma}-\bm{C}\rVert_{\mathrm{F}}^2.
\]
第二项非负。第一项在所有秩不超过$k$的正交投影中，由投影到前$k$个坐标方向达到
最小；其值为$\sum_{i>k}\sigma_i^2$。为看清这一步，注意
\[
  \lVert\bm{P}\bm{\Sigma}\rVert_{\mathrm{F}}^2
  =\sum_{i=1}^{q}\sigma_i^2\lVert\bm{P}\bm{e}_i\rVert_2^2,
  \qquad q=\min(m,n),
\]
其中$\bm{e}_i$是$\mathbb{R}^m$中的坐标向量，每个权重位于$[0,1]$。这些权重之和不超过
$\operatorname{tr}(\bm{P})\leq k$。要使保留能量最大，应把权重给最大的$k$个
$\sigma_i^2$。取$\bm{C}$为截断后的$\bm{\Sigma}$即可同时使第二项为零并达到下界。
\end{proof}

最佳近似未必唯一。对Frobenius范数，若$\sigma_k>\sigma_{k+1}$，截断SVD给出的
最佳矩阵唯一；若$\sigma_k=\sigma_{k+1}$，可以在这个重奇异值对应的子空间内选择
不同的$k$维保留方向，得到多个同样最优的矩阵。谱范数只约束最大的剩余误差，即使
边界处存在严格间隙，也可能允许其他最佳矩阵。定理保证两个误差最小值，却不能在所有
情形下唯一决定“应该保留哪一个方向”。

对$\bm{X}_{\widetilde{\varepsilon}}$取$k=1$，谱范数误差正是
$\sigma_{\min}\sim\widetilde{\varepsilon}/\sqrt{2}$。当数据精度或噪声尺度大于这个量时，把
矩阵视为秩一近似可能比试图区分两个近共线方向更可靠。但截断会引入偏差；若响应确实
包含沿弱方向的稳定信号，直接丢弃该方向也会损失信息。

\subsection{主成分分析在固定维数下保留最大能量}

设中心化数据矩阵$\bm{X}\in\mathbb{R}^{n\times d}$的每一行是一个样本。希望用
$k$维子空间表示数据，可以选择具有标准正交列的
$\bm{V}_k\in\mathbb{R}^{d\times k}$，并以
$\bm{X}\bm{V}_k\bm{V}_k^{\top}$重建数据。重建误差为
\[
  \lVert\bm{X}-\bm{X}\bm{V}_k\bm{V}_k^{\top}\rVert_{\mathrm{F}}^2.
\]
由勾股关系，这等于
\[
  \lVert\bm{X}\rVert_{\mathrm{F}}^2-
  \lVert\bm{X}\bm{V}_k\rVert_{\mathrm{F}}^2.
\]
若使用无偏样本协方差
$\bm{X}^{\top}\bm{X}/(n-1)$，投影后各方向的样本方差之和正是
$\lVert\bm{X}\bm{V}_k\rVert_{\mathrm{F}}^2/(n-1)$。所以最小化重建误差等价于
最大化保留方差。Eckart--Young--Mirsky定理表明，应取$\bm{X}$的前$k$个右奇异
向量。这就是
\term{主成分分析}（principal component analysis, PCA）的矩阵形式。

中心化条件不可省略。若不减去样本均值，第一主方向可能主要指向数据整体离原点的位置，
而不是样本围绕均值变化最大的方向。特征量纲也会影响结果：以米和毫米表示同一变量会
改变协方差矩阵及主成分，因此需要根据任务判断是否标准化。PCA找到的是线性子空间，
也不保证最大方差方向就是最有预测价值的方向。

\subsection{伪逆选择可解部分中的最小范数解}

对紧致SVD
$\bm{A}=\bm{U}_r\bm{\Sigma}_r\bm{V}_r^{\top}$，定义
\term{Moore--Penrose伪逆}（Moore--Penrose pseudoinverse）
\begin{equation}
  \bm{A}^{\dagger}
  =
  \bm{V}_r\bm{\Sigma}_r^{-1}\bm{U}_r^{\top}
  =
  \sum_{i=1}^{r}\frac{1}{\sigma_i}
  \bm{v}_i\bm{u}_i^{\top}.
  \label{eq:matrix-analysis-pseudoinverse}
\end{equation}
它把输出在$\operatorname{range}(\bm{A})$上的分量沿非零奇异方向反向缩放，并把
正交于值域的分量丢弃。

伪逆也可以不借助某个特定SVD来刻画：$\bm{A}^{\dagger}$是唯一满足
\begin{align*}
  \bm{A}\bm{A}^{\dagger}\bm{A}&=\bm{A},&
  \bm{A}^{\dagger}\bm{A}\bm{A}^{\dagger}&=\bm{A}^{\dagger},\\
  (\bm{A}\bm{A}^{\dagger})^{\top}
    &=\bm{A}\bm{A}^{\dagger},&
  (\bm{A}^{\dagger}\bm{A})^{\top}
    &=\bm{A}^{\dagger}\bm{A}
\end{align*}
的矩阵。这四个Moore--Penrose等式分别保证反演在可解部分上一致，并使输入、输出
两侧得到正交投影。下面分别验证存在性与唯一性。

由紧致SVD直接得到
\[
  \bm{A}\bm{A}^{\dagger}
  =\bm{U}_r\bm{U}_r^{\top}=:\bm{P}_{U},
  \qquad
  \bm{A}^{\dagger}\bm{A}
  =\bm{V}_r\bm{V}_r^{\top}=:\bm{P}_{V}.
\]
$\bm{P}_{U}$与$\bm{P}_{V}$都是对称正交投影，且
$\bm{P}_{U}\bm{A}=\bm{A}$、
$\bm{P}_{V}\bm{A}^{\dagger}=\bm{A}^{\dagger}$。于是逐项有
\[
  \bm{A}\bm{A}^{\dagger}\bm{A}=\bm{A},
  \qquad
  \bm{A}^{\dagger}\bm{A}\bm{A}^{\dagger}=\bm{A}^{\dagger},
  \qquad
  (\bm{A}\bm{A}^{\dagger})^{\top}=\bm{A}\bm{A}^{\dagger},
  \qquad
  (\bm{A}^{\dagger}\bm{A})^{\top}=\bm{A}^{\dagger}\bm{A}.
\]

再设$\bm{X}$是任意满足四个等式的矩阵。$\bm{A}\bm{X}$对称，并由
$\bm{A}\bm{X}\bm{A}=\bm{A}$可知它幂等；其值域包含于
$\operatorname{range}(\bm{A})$，同时对每个
$\bm{y}=\bm{A}\bm{z}$都有
$\bm{A}\bm{X}\bm{y}=\bm{A}\bm{z}=\bm{y}$。因此
$\bm{A}\bm{X}$是到$\operatorname{range}(\bm{A})$的正交投影，即
$\bm{A}\bm{X}=\bm{P}_{U}=\bm{A}\bm{A}^{\dagger}$。
类似地，$\bm{X}\bm{A}$对称且幂等，并且
\[
  \operatorname{null}(\bm{X}\bm{A})
  =\operatorname{null}(\bm{A}).
\]
若$\bm{X}\bm{A}\bm{z}=\bm{0}$，则
$\bm{A}\bm{z}=\bm{A}\bm{X}\bm{A}\bm{z}=\bm{0}$；反向包含关系可直接由
$\bm{A}\bm{z}=\bm{0}$得到。故
$\bm{X}\bm{A}$是到
$\operatorname{null}(\bm{A})^{\perp}
=\operatorname{range}(\bm{A}^{\top})$的正交投影，即
$\bm{X}\bm{A}=\bm{P}_{V}=\bm{A}^{\dagger}\bm{A}$。最后利用第二个Penrose等式，
\[
  \bm{X}
  =\bm{X}\bm{A}\bm{X}
  =\bm{X}\bm{A}\bm{A}^{\dagger}
  =\bm{A}^{\dagger}\bm{A}\bm{A}^{\dagger}
  =\bm{A}^{\dagger}.
\]
因此，四个Penrose等式唯一地确定伪逆。重奇异值对应的基即使发生旋转，
式~\eqref{eq:matrix-analysis-pseudoinverse}表示的线性映射仍不变。

\begin{theorem}[最小二乘与最小范数]
\label{thm:matrix-analysis-pseudoinverse-minimum-norm}
给定$\bm{A}\in\mathbb{R}^{m\times n}$与
$\bm{b}\in\mathbb{R}^m$，向量
$\bm{x}^{\dagger}=\bm{A}^{\dagger}\bm{b}$满足：
\begin{enumerate}
  \item 它最小化$\lVert\bm{A}\bm{x}-\bm{b}\rVert_2$；
  \item 在全部最小二乘解中，它的欧氏范数最小；
  \item 最小范数解是唯一的，但一般的最小二乘解未必唯一。
\end{enumerate}
\end{theorem}

\begin{proof}
把$\bm{U}$补成$\mathbb{R}^m$的标准正交基，并将
$\bm{b}$唯一分解为
\[
  \bm{b}=\bm{U}_r\bm{c}+\bm{b}_{\perp},
  \qquad
  \bm{U}_r^{\top}\bm{b}_{\perp}=\bm{0}.
\]
再把任意$\bm{x}$写成
$\bm{x}=\bm{V}_r\bm{z}+\bm{x}_0$，其中
$\bm{x}_0\in\operatorname{null}(\bm{A})$。正交性给出
\[
  \lVert\bm{A}\bm{x}-\bm{b}\rVert_2^2
  =\lVert\bm{\Sigma}_r\bm{z}-\bm{c}\rVert_2^2
  +\lVert\bm{b}_{\perp}\rVert_2^2.
\]
第二项不依赖$\bm{x}$，第一项在
$\bm{z}=\bm{\Sigma}_r^{-1}\bm{c}$时达到零。因此全部最小二乘解为
\[
  \bm{x}=
  \bm{V}_r\bm{\Sigma}_r^{-1}\bm{U}_r^{\top}\bm{b}
  +\bm{x}_0
  =\bm{A}^{\dagger}\bm{b}+\bm{x}_0.
\]
两项正交，所以
$\lVert\bm{x}\rVert_2^2
=\lVert\bm{A}^{\dagger}\bm{b}\rVert_2^2+\lVert\bm{x}_0\rVert_2^2$。
范数仅在$\bm{x}_0=\bm{0}$时最小，由此同时得到最小性与唯一性。
\end{proof}

若$\bm{A}$列满秩，则
$\bm{A}^{\dagger}=(\bm{A}^{\top}\bm{A})^{-1}\bm{A}^{\top}$；若行满秩，则
$\bm{A}^{\dagger}=\bm{A}^{\top}(\bm{A}\bm{A}^{\top})^{-1}$。这些公式适合推导，
却不意味着应在数值计算中显式形成逆矩阵。尤其是前一个公式会把条件数平方，后文会说明
它为何可能损失精度。

\section{矩阵范数与有效秩}
\label{sec:matrix-analysis-norms-effective-rank}

\subsection{不同范数回答不同的误差问题}

\term{矩阵范数}（matrix norm）把矩阵映射为非负标量，并满足正定性、齐次性和
三角不等式。若还满足
$\lVert\bm{A}\bm{B}\rVert\leq
\lVert\bm{A}\rVert\lVert\bm{B}\rVert$，就称为次乘性范数。范数不是矩阵“大小”的
唯一真值；选用哪一种，取决于希望控制逐元素误差、总体能量，还是最坏输入方向。

最常用的三类范数为
\begin{align}
  \lVert\bm{A}\rVert_{\mathrm{F}}
  &=\left(\sum_{i,j}a_{i,j}^2\right)^{1/2}
  =
  \left(\sum_i\sigma_i^2\right)^{1/2},
  \label{eq:matrix-analysis-frobenius-norm}\\
  \lVert\bm{A}\rVert_2
  &=\max_{\bm{x}\neq\bm{0}}
  \frac{\lVert\bm{A}\bm{x}\rVert_2}{\lVert\bm{x}\rVert_2}
  =
  \sigma_1,
  \label{eq:matrix-analysis-spectral-norm}\\
  \lVert\bm{A}\rVert_*
  &=\sum_i\sigma_i.
  \label{eq:matrix-analysis-nuclear-norm}
\end{align}
它们分别称为Frobenius范数、谱范数与核范数。Frobenius范数汇总全部元素或全部奇异
方向的平方能量；谱范数只看放大最强的单位输入；核范数把奇异值线性相加，常被用作
低秩的凸替代。

这三种范数属于同一个奇异值范数族。对$p\geq1$，定义
\[
  \lVert\bm{A}\rVert_{S_p}
  =\left(\sum_i\sigma_i^p\right)^{1/p},
\]
称为\term{Schatten $p$范数}（Schatten \(p\)-norm）。其中
$p=1$给出核范数，$p=2$给出Frobenius范数，而
$p\to\infty$的极限给出谱范数。若$0<p<1$，同一公式仍能强调小奇异值的稀疏性，
但不满足三角不等式，因而只是拟范数，不能称为凸范数。

谱范数由向量欧氏范数诱导，故
$\lVert\bm{A}\bm{x}\rVert_2\leq
\lVert\bm{A}\rVert_2\lVert\bm{x}\rVert_2$。Frobenius范数也次乘，并有
\begin{equation}
  \lVert\bm{A}\bm{B}\rVert_{\mathrm{F}}
  \leq\lVert\bm{A}\rVert_2\lVert\bm{B}\rVert_{\mathrm{F}}
  \leq
  \lVert\bm{A}\rVert_{\mathrm{F}}\lVert\bm{B}\rVert_{\mathrm{F}}.
  \label{eq:matrix-analysis-mixed-submultiplicativity}
\end{equation}
第一步可以逐列应用谱范数定义再求和；第二步使用
$\lVert\bm{A}\rVert_2\leq\lVert\bm{A}\rVert_{\mathrm{F}}$。

对秩为$r$的矩阵，奇异值表达式还给出
\begin{equation}
  \lVert\bm{A}\rVert_2\leq\lVert\bm{A}\rVert_{\mathrm{F}}\leq
  \sqrt{r}\lVert\bm{A}\rVert_2,
  \qquad
  \lVert\bm{A}\rVert_{\mathrm{F}}\leq\lVert\bm{A}\rVert_*\leq
  \sqrt{r}\lVert\bm{A}\rVert_{\mathrm{F}}.
  \label{eq:matrix-analysis-norm-rank-relations}
\end{equation}
维数或秩增大时，这些范数的比值可以随$\sqrt{r}$增长。因此“范数都等价”在有限维
空间中只表示它们诱导相同的收敛概念，不表示误差界中的常数可以忽略。

\subsection{诱导范数保留输入输出的语义}

对任意向量范数$\lVert\cdot\rVert$，相应的诱导矩阵范数定义为
\[
  \lVert\bm{A}\rVert
  =\max_{\bm{x}\neq\bm{0}}
  \frac{\lVert\bm{A}\bm{x}\rVert}{\lVert\bm{x}\rVert}.
\]
除谱范数外，常见的还有
\[
  \lVert\bm{A}\rVert_1=\max_j\sum_i|a_{i,j}|,
  \qquad
  \lVert\bm{A}\rVert_{\infty}=\max_i\sum_j|a_{i,j}|.
\]
前者是最大绝对列和，后者是最大绝对行和。若输入误差用$\ell_1$范数衡量，
$\lVert\bm{A}\rVert_1$控制最坏输出放大；若关心每个坐标的最大误差，则
$\lVert\bm{A}\rVert_{\infty}$更自然。

逐元素最大值$\max_{i,j}|a_{i,j}|$不是由向量$\ell_\infty$范数诱导的上述矩阵
范数；使用“最大范数”一词时必须写清定义。

\term{混合范数}（mixed norm）则先在一组坐标内聚合，再在组之间聚合。例如
\[
  \lVert\bm{A}\rVert_{2,1}
  =\sum_{i=1}^{m}
    \left(\sum_{j=1}^{n}a_{i,j}^2\right)^{1/2}
\]
先计算每一行的欧氏范数，再把各行结果相加，因而控制坐标分组而不只是奇异谱。也有
文献用$\lVert\cdot\rVert_{2,1}$表示按列分组的版本，所以使用时应同时说明分组方向。
这类范数依赖给定坐标中的行列分组，不像Schatten范数那样只由奇异值决定。

\subsection{有效秩描述奇异值能量是否集中}

代数秩把非零与零严格分开，却对数值精度极其敏感。任意微小的扰动都可能把零奇异值变成
非零值。若想描述一个矩阵“实际上有多少个重要方向”，需要连续的尺度。

对非零矩阵，一种常用定义是\term{稳定秩}（stable rank）
\begin{equation}
  r_{\mathrm{stable}}(\bm{A})
  =\frac{\lVert\bm{A}\rVert_{\mathrm{F}}^2}
        {\lVert\bm{A}\rVert_2^2}
  =
  \frac{\sum_i\sigma_i^2}{\sigma_1^2}.
  \label{eq:matrix-analysis-stable-rank}
\end{equation}
它位于$[1,\operatorname{rank}(\bm{A})]$内。所有非零奇异值相等时，它等于代数秩；
能量集中在第一个奇异值时，它接近$1$。
零矩阵会使定义式成为$0/0$，通常另约定
$r_{\mathrm{stable}}(\bm{0})=0$，而不是把区间$[1,\operatorname{rank}(\bm{A})]$
机械用于零矩阵。

稳定秩是常用的\term{有效秩}（effective rank）之一。文献中还有基于归一化奇异值
谱熵的有效秩，因此使用这个名称时必须同时给出公式。本节不把它称为“有效维数”：
第~\ref{chap:functional-analysis-and-approximation}章的核有效维度
$d_{\mathrm{eff}}(\lambda)$按正则化尺度$\lambda$给每个核谱方向加权，描述的是函数
空间在该尺度下的复杂度；稳定秩不含正则化尺度，比较的是一个有限矩阵的总平方奇异值
能量与最大平方奇异值。二者都能连续地计算谱方向的贡献，但对象和用途不同。

对$\bm{X}_{\widetilde{\varepsilon}}$，
\[
  r_{\mathrm{stable}}(\bm{X}_{\widetilde{\varepsilon}})
  =\frac{\sigma_{\max}^2+\sigma_{\min}^2}{\sigma_{\max}^2}
  =
  1+\frac{\lambda_-}{\lambda_+}
  \longrightarrow 1.
\]
它在$\widetilde{\varepsilon}>0$时代数秩始终为$2$，但稳定秩连续地反映出矩阵越来越接近秩一。
有效秩适合描述尺度，不会自动给出截断阈值；阈值还应结合噪声、任务损失与计算预算。

矩阵乘积的幅度可以逐层控制。若维度相容，则
\[
  \lVert\bm{A}_L\cdots\bm{A}_1\rVert_2
  \leq\prod_{\ell=1}^{L}\lVert\bm{A}_{\ell}\rVert_2.
\]
这给出深层线性映射或Jacobian乘积的最坏方向上界，但每一步取到最坏方向通常要求这些
方向能够对齐，所以上界未必紧。稳定秩也不能沿乘积简单相乘，更不等于截断以后实际
保留的整数方向数。

\section{矩阵函数与结构化运算}
\label{sec:matrix-analysis-functions-structured-operations}

\subsection{从标量函数到矩阵函数}

若对角化矩阵
$\bm{A}=\bm{V}\bm{\Lambda}\bm{V}^{-1}$的谱位于标量函数$f$的定义域内，可以定义
\begin{equation}
  f(\bm{A})=\bm{V}f(\bm{\Lambda})\bm{V}^{-1},
  \qquad
  f(\bm{\Lambda})=
  \operatorname{diag}\bigl(f(\lambda_1),\ldots,f(\lambda_d)\bigr).
  \label{eq:matrix-analysis-matrix-function-diagonalizable}
\end{equation}
白话地说，矩阵函数在每个特征方向上应用同一个标量函数。对称正定矩阵因而有唯一的
对称正定平方根
\[
  \bm{A}^{1/2}
  =\bm{Q}\operatorname{diag}(\sqrt{\lambda_i})\bm{Q}^{\top},
\]
并可定义
$\bm{A}^{-1/2}
=\bm{Q}\operatorname{diag}(\lambda_i^{-1/2})\bm{Q}^{\top}$、
$\log\bm{A}
=\bm{Q}\operatorname{diag}(\log\lambda_i)\bm{Q}^{\top}$与任意实数次幂。
严格正定保证这些表达都为实矩阵，也保证正定平方根唯一；半正定矩阵仍有唯一半正定
平方根，但零特征值使逆平方根和实对数不再按上述公式定义。

一般方阵未必可对角化。复Schur分解保证存在酉矩阵$\bm{Q}$和上三角矩阵$\bm{T}$使
\begin{equation}
  \bm{A}=\bm{Q}\bm{T}\bm{Q}^{*}.
  \label{eq:matrix-analysis-schur-decomposition}
\end{equation}
$\bm{T}$的对角元就是$\bm{A}$的特征值。Schur形式避免使用可能病态的特征向量矩阵，
许多数值矩阵函数算法会先得到$\bm{T}$，在三角矩阵上计算$f(\bm{T})$，再变换回来。
若只在实数中计算，也可以使用由$1\times1$与$2\times2$对角块组成的实Schur形式。

另一条定义路线来自幂级数。若
$f(z)=\sum_{k=0}^{\infty}a_kz^k$在包含矩阵谱的区域收敛，则
\[
  f(\bm{A})=\sum_{k=0}^{\infty}a_k\bm{A}^k.
\]
例如矩阵指数为
\begin{equation}
  \exp(\bm{A})=\sum_{k=0}^{\infty}\frac{\bm{A}^k}{k!}.
  \label{eq:matrix-analysis-matrix-exponential}
\end{equation}
级数对每个有限矩阵都收敛。若线性常微分方程满足
$\bm{x}'(t)=\bm{A}\bm{x}(t)$且$\bm{x}(0)=\bm{x}_0$，则
$\bm{x}(t)=\exp(t\bm{A})\bm{x}_0$。直接对级数逐项求导可验证这个结论。

幂级数定义也覆盖不可对角化矩阵。复数域上的Jordan分解把任意方阵相似变换为若干
Jordan块，每个块可写成
$\bm{J}=\lambda\bm{I}+\bm{N}$，其中$\bm{N}$只有超对角线为$1$，且对块大小
$s$满足$\bm{N}^{s}=\bm{0}$。因为$\lambda\bm{I}$与$\bm{N}$交换，
\[
  \exp(t\bm{J})
  =e^{\lambda t}
    \sum_{j=0}^{s-1}\frac{t^j}{j!}\bm{N}^j.
\]
求和在$s-1$次终止，说明不可对角化块除了指数因子，还会产生关于$t$的多项式因子。
这里使用Jordan形式是为了解释结构，不表示数值计算应显式求一个可能病态的Jordan基。
第~\ref{chap:dynamical-systems-and-state-estimation}章判断线性系统稳定性时，会同时
检查指数衰减和这些多项式因子。

矩阵乘法通常不交换，所以标量恒等式不能机械搬来。例如一般只有在
$\bm{A}\bm{B}=\bm{B}\bm{A}$时才有
\[
  \exp(\bm{A}+\bm{B})=\exp(\bm{A})\exp(\bm{B}).
\]
引入标量$t$并展开到二阶，可以直接看出差异：
\begin{align*}
  \exp(t(\bm{A}+\bm{B}))
  &=
  \bm{I}+t(\bm{A}+\bm{B})
  +\frac{t^2}{2}
   (\bm{A}^2+\bm{A}\bm{B}+\bm{B}\bm{A}+\bm{B}^2)
  +O(t^3),\\
  \exp(t\bm{A})\exp(t\bm{B})
  &=
  \bm{I}+t(\bm{A}+\bm{B})
  +t^2\left(
    \frac{\bm{A}^2}{2}+\bm{A}\bm{B}+\frac{\bm{B}^2}{2}
  \right)
  +O(t^3).
\end{align*}
左端的二阶交叉项为
$(\bm{A}\bm{B}+\bm{B}\bm{A})/2$，右端则为$\bm{A}\bm{B}$；两者之差为
$(\bm{B}\bm{A}-\bm{A}\bm{B})/2$。非交换性正是矩阵微积分和结构化算法中必须
持续检查的边界。

\subsection{分块消元产生Schur补}

分块矩阵并不只是记号上的拼接；消去一组变量以后，另一组变量会获得一个新的有效
系数。设
\[
  \bm{M}=
  \begin{bmatrix}
    \bm{A}&\bm{B}\\
    \bm{C}&\bm{D}
  \end{bmatrix},
\]
其中$\bm{A}$为可逆方阵。对块行做与高斯消元相同的操作，
\begin{equation}
  \begin{bmatrix}
    \bm{I}&\bm{0}\\
    -\bm{C}\bm{A}^{-1}&\bm{I}
  \end{bmatrix}
  \begin{bmatrix}
    \bm{A}&\bm{B}\\
    \bm{C}&\bm{D}
  \end{bmatrix}
  =
  \begin{bmatrix}
    \bm{A}&\bm{B}\\
    \bm{0}&\bm{S}
  \end{bmatrix},
  \qquad
  \bm{S}=\bm{D}-\bm{C}\bm{A}^{-1}\bm{B}.
  \label{eq:matrix-analysis-schur-complement}
\end{equation}
$\bm{S}$称为$\bm{A}$在$\bm{M}$中的\term{Schur补}（Schur complement）。
若右端按相同方式分块，先由第一组方程表达被消去的变量，再代入第二组方程，得到的正是
以$\bm{S}$为系数的约化方程。因此Schur补记录了消去第一组变量后，第二组变量之间剩余
的直接作用。

当$\bm{M}$对称且$\bm{C}=\bm{B}^{\top}$时，正定性也能沿这次消元传递。对任意
$\bm{x},\bm{y}$，配方得到
\[
  \begin{bmatrix}\bm{x}\\\bm{y}\end{bmatrix}^{\top}
  \bm{M}
  \begin{bmatrix}\bm{x}\\\bm{y}\end{bmatrix}
  =
  (\bm{x}+\bm{A}^{-1}\bm{B}\bm{y})^{\top}
  \bm{A}
  (\bm{x}+\bm{A}^{-1}\bm{B}\bm{y})
  +\bm{y}^{\top}\bm{S}\bm{y}.
\]
由此可见，在$\bm{A}\succ\bm{0}$的前提下，
\begin{equation}
  \bm{M}\succ\bm{0}
  \quad\Longleftrightarrow\quad
  \bm{S}=\bm{D}-\bm{B}^{\top}\bm{A}^{-1}\bm{B}
  \succ\bm{0}.
  \label{eq:matrix-analysis-schur-positive-definite}
\end{equation}
半正定情形不能只把两个严格不等号机械改成非严格不等号，因为$\bm{A}$可能不可逆；
此时需要伪逆以及相应的值域条件。

若$\bm{S}$也可逆，依次解两个三角分块系统可得
\begin{equation}
  \bm{M}^{-1}
  =
  \begin{bmatrix}
    \bm{A}^{-1}+\bm{A}^{-1}\bm{B}\bm{S}^{-1}
      \bm{C}\bm{A}^{-1}
    &-\bm{A}^{-1}\bm{B}\bm{S}^{-1}\\
    -\bm{S}^{-1}\bm{C}\bm{A}^{-1}
    &\bm{S}^{-1}
  \end{bmatrix}.
  \label{eq:matrix-analysis-block-inverse}
\end{equation}
式~\eqref{eq:matrix-analysis-schur-complement}至
\eqref{eq:matrix-analysis-block-inverse}将在第~\ref{chap:probability-theory}章
用于条件高斯的均值与协方差，并在
第~\ref{chap:dynamical-systems-and-state-estimation}章用于滤波和状态估计。
这些公式适合推导结构；实际计算仍应通过分解和线性求解应用
$\bm{A}^{-1}$与$\bm{S}^{-1}$，而不是显式形成逆矩阵。

\subsection{低秩更新不必重新求逆}

设$\bm{A}\in\mathbb{R}^{n\times n}$可逆，
$\bm{U},\bm{V}\in\mathbb{R}^{n\times k}$，
$\bm{C}\in\mathbb{R}^{k\times k}$可逆。若
$\bm{C}^{-1}+\bm{V}^{\top}\bm{A}^{-1}\bm{U}$也可逆，则
\term{Woodbury恒等式}给出
\begin{equation}
  (\bm{A}+\bm{U}\bm{C}\bm{V}^{\top})^{-1}
  =\bm{A}^{-1}-
  \bm{A}^{-1}\bm{U}
  \left(\bm{C}^{-1}+\bm{V}^{\top}\bm{A}^{-1}\bm{U}\right)^{-1}
  \bm{V}^{\top}\bm{A}^{-1}.
  \label{eq:matrix-analysis-woodbury}
\end{equation}
右侧把一个$n$维更新问题缩成$k$维求解。记括号内的小矩阵为$\bm{S}$，将右侧与
$\bm{A}+\bm{U}\bm{C}\bm{V}^{\top}$相乘，并代入
$\bm{I}-\bm{S}^{-1}\bm{V}^{\top}\bm{A}^{-1}\bm{U}
=\bm{S}^{-1}\bm{C}^{-1}$，即可验证乘积为单位矩阵。
它也可由分块矩阵
$\left[\begin{smallmatrix}\bm{A}&\bm{U}\\
-\bm{V}^{\top}&\bm{C}^{-1}\end{smallmatrix}\right]$
按两种顺序消元得到：消去$\bm{A}$产生小Schur补
$\bm{C}^{-1}+\bm{V}^{\top}\bm{A}^{-1}\bm{U}$，消去
$\bm{C}^{-1}$则产生更新后的大矩阵
$\bm{A}+\bm{U}\bm{C}\bm{V}^{\top}$。这说明Woodbury恒等式本质上是在比较同一个
分块系统的两种消元次序。

当$k=1$且$\bm{C}=1$时，该式退化为Sherman--Morrison公式。它常用于样本增删、
低秩协方差更新和带低秩修正的线性系统。但公式出现$\bm{A}^{-1}$不表示实现时应先形成
逆矩阵；更稳定的做法是复用$\bm{A}$的分解，通过线性求解得到
$\bm{A}^{-1}\bm{U}$。若中间的小矩阵接近奇异，更新后的大矩阵也接近奇异，公式本身
不会消除这个敏感性。

不含$\bm{C}$的常用形式更直接体现辅助变量消元。要解
$(\bm{A}+\bm{U}\bm{V}^{\top})\bm{x}=\bm{b}$，可先令
$\bm{z}=\bm{V}^{\top}\bm{x}$。由
$\bm{x}=\bm{A}^{-1}(\bm{b}-\bm{U}\bm{z})$得到
\[
  (\bm{I}+\bm{V}^{\top}\bm{A}^{-1}\bm{U})\bm{z}
  =\bm{V}^{\top}\bm{A}^{-1}\bm{b}.
\]
因此只需假设$\bm{A}$与
$\bm{I}+\bm{V}^{\top}\bm{A}^{-1}\bm{U}$可逆，不需要给通常为长方形的
$\bm{U}$或$\bm{V}$假定不存在的逆。

例如取
$\bm{A}=\operatorname{diag}(2,3)$、
$\bm{u}=\bm{v}=(1,1)^{\top}$和$\bm{b}=(1,0)^{\top}$。
直接求解
$\left[\begin{smallmatrix}3&1\\1&4\end{smallmatrix}\right]\bm{x}
=\bm{b}$得到
$\bm{x}=(4/11,-1/11)^{\top}$。辅助变量公式中的小系统只有一个标量，
\[
  1+\bm{u}^{\top}\bm{A}^{-1}\bm{u}
  =1+\frac12+\frac13=\frac{11}{6}.
\]
代入Sherman--Morrison公式也得到
\[
  \bm{x}
  =\bm{A}^{-1}\bm{b}
  -\bm{A}^{-1}\bm{u}
   \left(1+\bm{u}^{\top}\bm{A}^{-1}\bm{u}\right)^{-1}
   \bm{u}^{\top}\bm{A}^{-1}\bm{b}
  =\begin{bmatrix}4/11\\-1/11\end{bmatrix}.
\]
这个二阶算例展示的是计算组织方式；规模收益只有在更新秩远小于原矩阵维数且原分解
可以复用时才出现。

\subsection{Kronecker积表达可分离结构}

沿用式~\eqref{eq:linear-algebra-kronecker-product}中的Kronecker积定义。给定
$\bm{A}\in\mathbb{R}^{m\times n}$与
$\bm{B}\in\mathbb{R}^{p\times q}$，
$\bm{A}\otimes\bm{B}$的维度是$mp\times nq$。只要各乘积维度相容，就有
\begin{equation}
  (\bm{A}\otimes\bm{B})(\bm{C}\otimes\bm{D})
  =(\bm{A}\bm{C})\otimes(\bm{B}\bm{D}).
  \label{eq:matrix-analysis-kronecker-mixed-product}
\end{equation}
若$\bm{A}$与$\bm{B}$均为可逆方阵，还可得到
$(\bm{A}\otimes\bm{B})^{-1}
=\bm{A}^{-1}\otimes\bm{B}^{-1}$。

把矩阵按列堆叠为向量的运算记为$\operatorname{vec}$。逐列检查可得
\begin{equation}
  \operatorname{vec}(\bm{A}\bm{X}\bm{B})
  =(\bm{B}^{\top}\otimes\bm{A})\operatorname{vec}(\bm{X}).
  \label{eq:matrix-analysis-vec-kronecker}
\end{equation}
若$\bm{A}\in\mathbb{R}^{p\times m}$、
$\bm{X}\in\mathbb{R}^{m\times n}$、
$\bm{B}\in\mathbb{R}^{n\times q}$，则等式左侧和右侧都在
$\mathbb{R}^{pq}$中。更具体地，$\bm{A}\bm{X}\bm{B}$的第$(i,j)$个元素为
\[
  \sum_{\ell=1}^{m}\sum_{k=1}^{n}
  a_{i,\ell}x_{\ell,k}b_{k,j}.
\]
按列堆叠后，这个元素位于索引$i+(j-1)p$；展开
$(\bm{B}^{\top}\otimes\bm{A})\operatorname{vec}(\bm{X})$在同一索引的分量，
得到完全相同的二重求和。这既证明了恒等式，也解释了为什么右侧是
$\bm{B}^{\top}\otimes\bm{A}$而不是两个因子的其他次序。
这个恒等式把矩阵方程转成线性方程，也揭示了多维网格与可分离协方差中的结构。实际计算
通常不应显式构造巨大的Kronecker矩阵；利用左右乘法实现同一个线性算子，可以显著减少
存储和运算量。

\subsection{极分解分开旋转与伸缩}

任意$\bm{A}\in\mathbb{R}^{m\times n}$都有
\term{极分解}（polar decomposition）
\begin{equation}
  \bm{A}=\bm{Q}\bm{H},
  \qquad
  \bm{H}=(\bm{A}^{\top}\bm{A})^{1/2},
  \label{eq:matrix-analysis-polar-decomposition}
\end{equation}
其中$\bm{H}$半正定，$\bm{Q}$是部分等距映射。用SVD可取
$\bm{H}=\bm{V}(\bm{\Sigma}^{\top}\bm{\Sigma})^{1/2}\bm{V}^{\top}$；
更直接地，在紧致SVD下，
$\bm{H}=\bm{V}_r\bm{\Sigma}_r\bm{V}_r^{\top}$，并在零空间上取零，
$\bm{Q}=\bm{U}_r\bm{V}_r^{\top}\in\mathbb{R}^{m\times n}$。此时
\[
  \bm{Q}^{\top}\bm{Q}
  =\bm{V}_r\bm{V}_r^{\top}
  =\bm{P}_{\operatorname{null}(\bm{A})^{\perp}},
  \qquad
  \bm{Q}\bm{Q}^{\top}
  =\bm{U}_r\bm{U}_r^{\top}
  =\bm{P}_{\operatorname{range}(\bm{A})}.
\]
因此“部分等距”是指$\bm{Q}$在
$\operatorname{null}(\bm{A})^{\perp}$上保持欧氏长度，并把该子空间等距映到
$\operatorname{range}(\bm{A})$；它把$\operatorname{null}(\bm{A})$中的方向映为零。

当$\bm{A}$为可逆方阵时，$\bm{Q}$正交且分解唯一。矩阵秩亏时，$\bm{H}$仍唯一，
要求$\operatorname{null}(\bm{Q})=\operatorname{null}(\bm{A})$时，上式给出的
部分等距因子也唯一；若维数允许，并进一步把它延拓为在零空间上仍保持长度的等距映射，
则原矩阵没有规定这些零方向应映到哪里，延拓可能不唯一。极分解把矩阵的
角度变化与非负伸缩分开，在正交约束优化、形状对齐和迭代矩阵正交化中很有用。

\subsection{矩阵函数的Fréchet导数保留非交换信息}

第~\ref{chap:mathematical-analysis}章已经从乘积恒等式推导了逆矩阵的微分，并用
Jacobi公式推导了$\log\det\bm{A}$的微分。本节关心更一般的问题：当输入矩阵沿
$\bm{E}$变化时，矩阵函数$f(\bm{A})$的一阶变化是什么。若存在线性映射
$L_f(\bm{A},\cdot)$使
\begin{equation}
  f(\bm{A}+\bm{E})-f(\bm{A})
  =
  L_f(\bm{A},\bm{E})+o(\lVert\bm{E}\rVert),
  \label{eq:matrix-analysis-frechet-derivative}
\end{equation}
就称$L_f(\bm{A},\bm{E})$为$f$在$\bm{A}$处沿$\bm{E}$的
\term{Fréchet导数}（Fréchet derivative）。它对扰动$\bm{E}$线性，却通常不是
简单的$f'(\bm{A})\bm{E}$。

对称矩阵$\bm{A}=\bm{Q}\operatorname{diag}(\lambda_i)\bm{Q}^{\top}$及对称扰动
$\bm{E}$给出一个直接表达。定义一阶差商矩阵
\[
  \bm{F}^{[1]}_{i,j}
  =
  \begin{cases}
    \dfrac{f(\lambda_i)-f(\lambda_j)}{\lambda_i-\lambda_j},
      &\lambda_i\neq\lambda_j,\\[6pt]
    f'(\lambda_i),&\lambda_i=\lambda_j.
  \end{cases}
\]
在$f$对相应谱区间连续可微等保证矩阵函数可微的条件下，
\begin{equation}
  L_f(\bm{A},\bm{E})
  =
  \bm{Q}\left(
    \bm{F}^{[1]}\mathbin{\odot}
    (\bm{Q}^{\top}\bm{E}\bm{Q})
  \right)\bm{Q}^{\top},
  \label{eq:matrix-analysis-frechet-divided-difference}
\end{equation}
其中$\odot$表示逐元素乘法。对角项使用标量导数，非对角项却由两个不同特征值之间的
差商控制；这正是标量口诀遗漏的非交换信息。当
$\bm{A}\bm{E}=\bm{E}\bm{A}$时，扰动在$\bm{A}$的不同谱块之间没有交叉，
式~\eqref{eq:matrix-analysis-frechet-divided-difference}可以化为熟悉的
$L_f(\bm{A},\bm{E})=f'(\bm{A})\bm{E}$。若二者不交换，一般不能作此简化；
不过对仿射函数等特殊情形，等式仍可能成立，所以交换性是这里的充分条件而非必要条件。

矩阵指数把这条边界表现得很清楚：
\begin{equation}
  L_{\exp}(\bm{A},\bm{E})
  =
  \int_0^1
  \exp((1-s)\bm{A})\bm{E}\exp(s\bm{A})\,\mathrm{d}s.
  \label{eq:matrix-analysis-exponential-frechet-derivative}
\end{equation}
若$\bm{A}$与$\bm{E}$交换，积分才等于
$\exp(\bm{A})\bm{E}$；一般情形必须保留$\bm{E}$两侧因子的次序。逆矩阵微分中的
$-\bm{A}^{-1}\bm{E}\bm{A}^{-1}$是同一现象的特例，而标量函数
$\log\det\bm{A}$输出一个数，其迹表达式已由第~\ref{chap:mathematical-analysis}章
处理，无须在这里重复推导。

平方根的微分给出另一个不宜交换次序的例子。设$\bm{A}\succ\bm{0}$且
$\bm{S}=\bm{A}^{1/2}$。沿对称扰动$\bm{E}$对恒等式
$\bm{S}^2=\bm{A}$求一阶变化，若
$\bm{Z}=L_{\sqrt{\cdot}}(\bm{A},\bm{E})$，则$\bm{Z}$满足
\[
  \bm{S}\bm{Z}+\bm{Z}\bm{S}=\bm{E}.
\]
这是一类\term{Sylvester方程}（Sylvester equation）。在
$\bm{S}=\bm{Q}\operatorname{diag}(s_i)\bm{Q}^{\top}$的特征基中，记
$\widehat{\bm{Z}}=\bm{Q}^{\top}\bm{Z}\bm{Q}$与
$\widehat{\bm{E}}=\bm{Q}^{\top}\bm{E}\bm{Q}$，逐项便有
\[
  \widehat{z}_{i,j}
  =\frac{\widehat{e}_{i,j}}{s_i+s_j}.
\]
正定性保证$s_i+s_j>0$，所以解唯一。只有在$\bm{A}$与$\bm{E}$交换时，非对角
耦合消失，结果才可简化为标量式
$\bm{Z}=\tfrac12\bm{A}^{-1/2}\bm{E}$；一般情形中把
$\bm{A}^{-1/2}$随意放到扰动一侧都会丢失次序信息。

\section{条件数与扰动解释何时难以求准}
\label{sec:matrix-analysis-conditioning-perturbation}

\subsection{条件数属于问题，而不是某个算法}

对可逆方阵$\bm{A}$，相对于诱导范数的
\term{条件数}（condition number）定义为
\begin{equation}
  \kappa(\bm{A})
  =\lVert\bm{A}\rVert\lVert\bm{A}^{-1}\rVert.
  \label{eq:matrix-analysis-condition-number}
\end{equation}
在谱范数下，
\begin{equation}
  \kappa_2(\bm{A})
  =\frac{\sigma_{\max}(\bm{A})}{\sigma_{\min}(\bm{A})}.
  \label{eq:matrix-analysis-spectral-condition-number}
\end{equation}
它比较最强与最弱的方向性缩放。正交矩阵的条件数为$1$；矩阵接近奇异时，
$\sigma_{\min}$接近零，条件数变大。秩亏矩阵通常约定条件数为无穷大。
对满秩长方矩阵$\bm{A}\in\mathbb{R}^{m\times n}$，常把谱条件数定义为
$\kappa_2(\bm{A})=\sigma_1/\sigma_q$，其中$q=\min(m,n)$且$\sigma_q>0$。
若矩阵秩亏，涉及反演或最小二乘问题在任意输入扰动下的条件数仍通常约定为无穷大；
但固定该矩阵以后，伪逆作为线性映射仍有有限范数，这两种口径需要区分。

只扰动右端时，令$\bm{b}\neq\bm{0}$，并考虑
$\bm{A}\bm{x}=\bm{b}$和
$\bm{A}(\bm{x}+\Delta\bm{x})=\bm{b}+\Delta\bm{b}$。两式相减得到
$\Delta\bm{x}=\bm{A}^{-1}\Delta\bm{b}$，于是
\[
  \frac{\lVert\Delta\bm{x}\rVert}
       {\lVert\bm{x}\rVert}
  \leq
  \kappa(\bm{A})
  \frac{\lVert\Delta\bm{b}\rVert}
       {\lVert\bm{b}\rVert}.
\]
这里用到了
$\lVert\bm{b}\rVert=\lVert\bm{A}\bm{x}\rVert
\leq\lVert\bm{A}\rVert\lVert\bm{x}\rVert$。该不等式是最坏方向的上界，并不表示每个
扰动都会被放大$\kappa(\bm{A})$倍。

对$\bm{X}_{\widetilde{\varepsilon}}$，
\[
  \kappa_2(\bm{X}_{\widetilde{\varepsilon}})
  =\sqrt{\frac{\lambda_+}{\lambda_-}}
  \sim\frac{2}{\widetilde{\varepsilon}}.
\]
式~\eqref{eq:matrix-analysis-collinear-coefficient-perturbation}选择的
$\Delta\bm{y}=(0,\eta)^{\top}$含有弱左奇异方向上的分量，因此系数变化达到
$1/\widetilde{\varepsilon}$量级。这不是求解器制造的不稳定，而是从响应反推两个近共线系数本来
就敏感。

\subsection{同时扰动矩阵与右端}

设$\bm{A}\in\mathbb{R}^{d\times d}$可逆、$\bm{b}\neq\bm{0}$，原系统为
$\bm{A}\bm{x}=\bm{b}$，其唯一解为$\bm{x}$。扰动后为
\[
  (\bm{A}+\Delta\bm{A})(\bm{x}+\Delta\bm{x})
  =\bm{b}+\Delta\bm{b}.
\]
整理可得
\[
  \bm{A}\Delta\bm{x}
  =\Delta\bm{b}-\Delta\bm{A}\bm{x}-\Delta\bm{A}\Delta\bm{x}.
\]
若
$\lVert\bm{A}^{-1}\rVert\lVert\Delta\bm{A}\rVert<1$，把最后一项移到左侧并
用范数界控制，得到
\begin{equation}
  \frac{\lVert\Delta\bm{x}\rVert}{\lVert\bm{x}\rVert}
  \leq
  \frac{\kappa(\bm{A})}
       {1-\kappa(\bm{A})
       \lVert\Delta\bm{A}\rVert/\lVert\bm{A}\rVert}
  \left(
    \frac{\lVert\Delta\bm{A}\rVert}{\lVert\bm{A}\rVert}
    +
    \frac{\lVert\Delta\bm{b}\rVert}{\lVert\bm{b}\rVert}
  \right).
  \label{eq:matrix-analysis-linear-system-perturbation}
\end{equation}
分母条件保证扰动还没有把系统推过奇异边界。若条件不满足，这个上界失效，并不等于
扰动后的矩阵一定奇异；它只表示该推导不能再保证稳定。

正规方程会恶化近共线问题。列满秩矩阵满足
\begin{equation}
  \kappa_2(\bm{A}^{\top}\bm{A})=\kappa_2(\bm{A})^2,
  \label{eq:matrix-analysis-normal-equation-condition-square}
\end{equation}
因为$\bm{A}^{\top}\bm{A}$的最大、最小特征值分别是
$\sigma_{\max}^2$与$\sigma_{\min}^2$。对$\bm{X}_{\widetilde{\varepsilon}}$，条件数从
$2/\widetilde{\varepsilon}$量级变为$4/\widetilde{\varepsilon}^2$量级。这就是数值计算中更偏好QR或SVD
直接处理最小二乘问题的原因。

\subsection{系数不稳定不等于拟合值不稳定}

设$\bm{X}\in\mathbb{R}^{n\times d}$，最小二乘拟合值为
\[
  \widehat{\bm{y}}=\bm{X}\bm{X}^{\dagger}\bm{y}.
\]
矩阵$\bm{P}_{\bm{X}}=\bm{X}\bm{X}^{\dagger}$是到
$\operatorname{range}(\bm{X})$的正交投影，所以
$\lVert\bm{P}_{\bm{X}}\rVert_2\leq1$。当设计矩阵固定而只扰动响应时，
\[
  \lVert\Delta\widehat{\bm{y}}\rVert_2
  =\lVert\bm{P}_{\bm{X}}\Delta\bm{y}\rVert_2
  \leq
  \lVert\Delta\bm{y}\rVert_2.
\]
相反，最小范数系数的变化
$\Delta\widehat{\bm{\beta}}=\bm{X}^{\dagger}\Delta\bm{y}$由
$\lVert\bm{X}^{\dagger}\rVert_2$控制。若
$r=\operatorname{rank}(\bm{X})>0$，则
\[
  \lVert\bm{X}^{\dagger}\rVert_2=\frac{1}{\sigma_r},
\]
其中$\sigma_r$是最小正奇异值；零奇异方向由伪逆置零，并不会按$1/0$放大。
若$\bm{X}=\bm{0}$，则$\bm{X}^{\dagger}=\bm{0}$。

这解释了贯穿例中的反差：两个大幅变化的系数沿近似相消方向移动，乘回
$\bm{X}_{\widetilde{\varepsilon}}$后只留下很小的预测变化。若任务目标是区分两个特征各自的作用，
这种敏感性是严重问题；若只关心训练设计空间内的预测，系数不稳定未必带来同等程度的
预测不稳定。对新样本的外推仍可能敏感，因为新样本未必沿训练数据中相同的近共线关系
变化。

岭正则化用
\[
  \widehat{\bm{\beta}}_{\lambda}
  =(\bm{X}^{\top}\bm{X}+\lambda\bm{I})^{-1}\bm{X}^{\top}\bm{y},
  \qquad \lambda>0,
\]
把每个$\sigma_i>0$的右奇异方向上的逆缩放从$1/\sigma_i$改为
\begin{equation}
  \frac{\sigma_i}{\sigma_i^2+\lambda}.
  \label{eq:matrix-analysis-ridge-filter-factor}
\end{equation}
零奇异方向在伪逆与岭映射中都给出零系数；正奇异值很小的方向不再被大幅放大，但估计
会向零收缩。正则化是在偏差与敏感性之间作选择，不是把原问题无代价地“修好”。
$\lambda$的尺度也依赖特征缩放，不能脱离数据单位比较。

\subsection{特征值与子空间对扰动的反应不同}

对称矩阵的特征值相对稳定。若$\bm{A}$与
$\widetilde{\bm{A}}=\bm{A}+\bm{E}$都对称，Weyl不等式给出
\begin{equation}
  |\widetilde{\lambda}_i-\lambda_i|
  \leq\lVert\bm{E}\rVert_2,
  \label{eq:matrix-analysis-weyl-eigenvalue-perturbation}
\end{equation}
其中两组特征值都按降序排列\scite{math-horn2013matrix}

这个界可以由定理~\ref{thm:matrix-analysis-courant-fischer}逐指标证明。令
$E_i=\operatorname{span}\{\bm{q}_1,\ldots,\bm{q}_i\}$、
$F_i=\operatorname{span}\{\bm{q}_i,\ldots,\bm{q}_d\}$，其中
$\bm{q}_j$是$\bm{A}$的第$j$个标准正交特征向量。对任意单位向量$\bm{x}$，
\[
  -\lVert\bm{E}\rVert_2
  \leq\bm{x}^{\top}\bm{E}\bm{x}
  \leq\lVert\bm{E}\rVert_2.
\]
在Courant--Fischer公式的第一个表达式中选$S=E_i$，得到
\[
  \widetilde{\lambda}_i
  \geq
  \min_{\substack{\bm{x}\in E_i\\\lVert\bm{x}\rVert_2=1}}
  \bm{x}^{\top}(\bm{A}+\bm{E})\bm{x}
  \geq\lambda_i-\lVert\bm{E}\rVert_2.
\]
在第二个表达式中选$T=F_i$，则
\[
  \widetilde{\lambda}_i
  \leq
  \max_{\substack{\bm{x}\in F_i\\\lVert\bm{x}\rVert_2=1}}
  \bm{x}^{\top}(\bm{A}+\bm{E})\bm{x}
  \leq\lambda_i+\lVert\bm{E}\rVert_2.
\]
上下界合并便得到式~\eqref{eq:matrix-analysis-weyl-eigenvalue-perturbation}，而且论证对
每个$i$分别成立。

比较特征方向和奇异方向之前，先统一子空间夹角的记号。设
$\bm{Q},\widetilde{\bm{Q}}\in\mathbb{R}^{d\times r}$都具有标准正交列，它们
张成的两个$r$维子空间之间的\term{主角}（principal angles）
$0\leq\theta_1\leq\cdots\leq\theta_r\leq\pi/2$由
$\bm{Q}^{\top}\widetilde{\bm{Q}}$的奇异值
$\cos\theta_1,\ldots,\cos\theta_r$定义。记
\[
  \bm{\Theta}(\bm{Q},\widetilde{\bm{Q}})
  =
  \operatorname{diag}(\theta_1,\ldots,\theta_r),
  \qquad
  \sin\bm{\Theta}(\bm{Q},\widetilde{\bm{Q}})
  =
  \operatorname{diag}(\sin\theta_1,\ldots,\sin\theta_r).
\]
这些量只依赖子空间，不依赖其中选择的标准正交基。一维时，
$\angle(\bm{q},\widetilde{\bm{q}})
=\arccos|\bm{q}^{\top}\widetilde{\bm{q}}|$就是唯一的主角；绝对值消除了特征向量
正负号的任意性。若$\bm{P}=\bm{Q}\bm{Q}^{\top}$且
$\widetilde{\bm{P}}=\widetilde{\bm{Q}}\widetilde{\bm{Q}}^{\top}$，则
\[
  \lVert\bm{P}-\widetilde{\bm{P}}\rVert_2
  =
  \lVert\sin\bm{\Theta}(\bm{Q},\widetilde{\bm{Q}})\rVert_2
  =
  \sin\theta_r.
\]

一个二阶算例先展示谱隙为何出现在方向误差的分母中。令
\[
  \bm{A}_{\gamma}
  =
  \begin{bmatrix}1+\gamma&0\\0&1\end{bmatrix},
  \qquad
  \bm{E}_{\epsilon}
  =
  \begin{bmatrix}0&\epsilon\\\epsilon&0\end{bmatrix},
  \qquad \gamma>0,\quad \epsilon>0.
\]
原矩阵最大特征值与其余特征值的间隙是$\gamma$，扰动范数为$\epsilon$。若扰动后
最大特征向量与$\bm{e}_1$的夹角为$\theta\in(0,\pi/4)$，直接解二阶对称矩阵的
特征方程可得
\[
  \tan(2\theta)=\frac{2\epsilon}{\gamma}.
\]
当$\epsilon\ll\gamma$时，$\theta\approx\epsilon/\gamma$。因此，在扰动大小
相同的情况下，谱隙越小，特征方向越容易旋转；当扰动与谱隙同量级时，线性近似也不再
足以描述转角。

特征向量却依赖谱间隙。设$\lambda_i$是单特征值，
$\bm{q}_i$为对应单位特征向量，扰动后的对应方向为
$\widetilde{\bm{q}}_i$，对应特征值为$\widetilde{\lambda}_i$。定义它与原矩阵其余
特征值的实际分离量
\[
  \operatorname{sep}_i
  =
  \min_{j\neq i}|\widetilde{\lambda}_i-\lambda_j|.
\]
若$\operatorname{sep}_i>0$，一维Davis--Kahan残差界给出
\begin{equation}
  \sin\angle(\bm{q}_i,\widetilde{\bm{q}}_i)
  \leq
  \frac{\lVert\bm{E}\rVert_2}{\operatorname{sep}_i}.
  \label{eq:matrix-analysis-eigenvector-gap-bound}
\end{equation}
这是把扰动特征方程投影到$\bm{q}_i^{\perp}$后，对
$\bm{A}-\widetilde{\lambda}_i\bm{I}$在该子空间上求逆得到的
\scite{math-horn2013matrix}若记原始间隙
$\delta_i=\min_{j\neq i}|\lambda_i-\lambda_j|$，Weyl不等式还给出
$\operatorname{sep}_i\geq\delta_i-\lVert\bm{E}\rVert_2$。这个形式明确显示：
真正控制方向变化的是扰动后目标值与原补谱之间仍然剩下的间隙。

多个相邻特征值组成一簇时，应直接比较子空间。若
$\bm{Q}_r$与$\widetilde{\bm{Q}}_r$的列分别张成两个$r$维子空间，记它们的正交
投影为
$\bm{P}=\bm{Q}_r\bm{Q}_r^{\top}$与
$\widetilde{\bm{P}}
=\widetilde{\bm{Q}}_r\widetilde{\bm{Q}}_r^{\top}$，定义距离
\[
  d(\bm{Q}_r,\widetilde{\bm{Q}}_r)
  =\lVert\bm{P}-\widetilde{\bm{P}}\rVert_2
  =\lVert\sin\bm{\Theta}(\bm{Q}_r,\widetilde{\bm{Q}}_r)\rVert_2.
\]
这个量不随子空间内部选择哪一组标准正交基而改变。若
$\operatorname{sep}_r=\widetilde{\lambda}_r-\lambda_{r+1}>0$，则前$r$维
特征子空间满足Davis--Kahan界
\[
  \lVert\bm{P}-\widetilde{\bm{P}}\rVert_2
  \leq\frac{\lVert\bm{E}\rVert_2}{\operatorname{sep}_r}.
\]
证明的关键不是逐个匹配簇内特征向量，而是把扰动后的特征方程投影到原子空间的正交
补；该补空间上的谱与扰动后目标谱簇至少相隔$\operatorname{sep}_r$，因而可以除以
这个间隙来控制泄漏分量\scite{math-horn2013matrix}若原边界间隙
$\delta=\lambda_r-\lambda_{r+1}>0$，则Weyl不等式给出
$\operatorname{sep}_r\geq\delta-\lVert\bm{E}\rVert_2$；只有这个余量为正时，
上述前$r$维谱簇才保持分离。

重特征值揭示了这个结论的边界。若
$\bm{A}=\lambda\bm{I}$，任意单位向量都是特征向量；一个任意小的对称扰动都可以选出
完全不同的坐标方向。此时要求某个特征向量稳定没有意义，稳定对象应改为整个重特征值
对应的不变子空间。类似地，PCA中若第$k$与第$k+1$个奇异值非常接近，前$k$维子空间
的边界容易变化；若关心的是包含整个重数块的更大子空间，则可能仍然稳定。
例如$\bm{A}=\operatorname{diag}(2,2,0)$的前二维特征空间可用
$\bm{e}_1,\bm{e}_2$表示，也可用任意角度$\theta$旋转后的
$\cos\theta\,\bm{e}_1+\sin\theta\,\bm{e}_2$与
$-\sin\theta\,\bm{e}_1+\cos\theta\,\bm{e}_2$表示。基向量可旋转任意角度，
对应投影却始终是$\operatorname{diag}(1,1,0)$。

\subsection{奇异向量与奇异子空间也需要谱隙}

一般矩阵的奇异值可通过对称扩张接入前述扰动结论。对
$\bm{A}\in\mathbb{R}^{m\times n}$定义
\[
  \mathcal{H}(\bm{A})
  =
  \begin{bmatrix}
    \bm{0}&\bm{A}\\
    \bm{A}^{\top}&\bm{0}
  \end{bmatrix}.
\]
若$(\bm{u}_i,\bm{v}_i,\sigma_i)$是一个奇异三元组，则
$2^{-1/2}(\bm{u}_i^{\top},\bm{v}_i^{\top})^{\top}$与
$2^{-1/2}(\bm{u}_i^{\top},-\bm{v}_i^{\top})^{\top}$分别是
$\mathcal{H}(\bm{A})$对应于$\sigma_i$与$-\sigma_i$的特征向量。扰动
$\widetilde{\bm{A}}=\bm{A}+\bm{E}$在扩张后仍是对称扰动，且
$\lVert\mathcal{H}(\bm{E})\rVert_2=\lVert\bm{E}\rVert_2$。Weyl不等式于是给出
\begin{equation}
  |\widetilde{\sigma}_i-\sigma_i|
  \leq\lVert\bm{E}\rVert_2.
  \label{eq:matrix-analysis-singular-value-perturbation}
\end{equation}
奇异值本身不需要谱隙也能稳定变化，但要识别对应方向，还必须把目标奇异值与其余谱
分开\scite{math-horn2013matrix}

具体地，设$\bm{U}_k,\bm{V}_k$分别张成$\bm{A}$的前$k$个左、右奇异子空间，
$\widetilde{\bm{U}}_k,\widetilde{\bm{V}}_k$是扰动矩阵的对应子空间，并令边界谱隙
\[
  \operatorname{sep}_{\mathrm{W}}
  =\widetilde{\sigma}_k-\sigma_{k+1}>0,
\]
其中越过矩阵秩的奇异值按零计。令
\[
  \bm{R}
  =\bm{A}\widetilde{\bm{V}}_k
   -\widetilde{\bm{U}}_k\widetilde{\bm{\Sigma}}_k,
  \qquad
  \bm{S}
  =\bm{A}^{\top}\widetilde{\bm{U}}_k
   -\widetilde{\bm{V}}_k\widetilde{\bm{\Sigma}}_k.
\]
由于$\widetilde{\bm{A}}=\bm{A}+\bm{E}$，这两个残差分别等于
$-\bm{E}\widetilde{\bm{V}}_k$与
$-\bm{E}^{\top}\widetilde{\bm{U}}_k$。Wedin正弦定理的Frobenius范数版本给出
\begin{equation}
  \left(
    \lVert\sin\bm{\Theta}(\bm{U}_k,\widetilde{\bm{U}}_k)\rVert_{\mathrm{F}}^2
    +
    \lVert\sin\bm{\Theta}(\bm{V}_k,\widetilde{\bm{V}}_k)\rVert_{\mathrm{F}}^2
  \right)^{1/2}
  \leq
  \frac{
    \left(\lVert\bm{R}\rVert_{\mathrm{F}}^2
    +\lVert\bm{S}\rVert_{\mathrm{F}}^2\right)^{1/2}
  }{\operatorname{sep}_{\mathrm{W}}}.
  \label{eq:matrix-analysis-singular-subspace-gap}
\end{equation}
左端使用前面定义的主角，同时汇总左右奇异子空间的偏离。这个残差形式避免在未固定范数
与分离量定义时写入额外常数，并直接显示子空间误差由“奇异方程没有满足到什么程度”与
“目标簇还剩多大间隙”共同控制\scite{math-golub2013matrix}
这个结论允许前$k$个奇异值内部重复：只要第$k$与第$k+1$个奇异值之间存在间隙，
整个前$k$维子空间仍可稳定；若还要辨认其中某个单独奇异向量，则该奇异值与两侧相邻
奇异值都要分离。没有边界谱隙时，小扰动可以在几乎等价的方向之间大幅旋转，此时只能
讨论包含完整奇异值簇的子空间，不能把某个数值算法返回的基向量当作可辨识对象。

\section{稳定求解把结构转化为可靠计算}
\label{sec:matrix-analysis-stable-solvers}

评价下面的求解方法时，“稳定”主要指\term{后向稳定}（backward stability）：
算法返回的结果是某个邻近输入问题的精确解。它衡量算法引入的等效输入扰动，不能单独
保证计算结果接近原问题的真解；后者还取决于问题的条件数。本节末尾将用残差、前向误差
与后向误差把这一区别形式化。

\subsection{消元与分解避免重复工作}

求解$\bm{A}\bm{x}=\bm{b}$时，显式计算$\bm{A}^{-1}$再乘$\bm{b}$通常不是合适的
实现。高斯消元把$\bm{A}$化为三角形式，随后用前代与回代求解。若有多个右端，同一个
分解可以重复使用。对一般稠密方阵，带主元选取的LU分解写作
\begin{equation}
  \bm{P}\bm{A}=\bm{L}\bm{U},
  \label{eq:matrix-analysis-palu}
\end{equation}
其中$\bm{P}$是置换矩阵，$\bm{L}$为单位下三角矩阵，$\bm{U}$为上三角矩阵。
对非奇异矩阵，主元选取可避免因当前候选主元为零而中断，并抑制消元过程中元素的过度
增长；秩亏矩阵仍会出现零主元，此时算法应报告秩亏或改用相应的分解。

主元选取不能改善问题的条件数，只抑制计算过程额外放大舍入误差。矩阵具有对称、正定、
带状或稀疏结构时，还应选用保留相应结构的分解。

这种说法依赖浮点计算的基本模型。排除上溢、下溢等异常后，一次基本运算通常写成
\[
  \operatorname{fl}(a\mathbin{\circ}b)
  =(a\mathbin{\circ}b)(1+\delta),
  \qquad |\delta|\leq u,
\]
其中$\circ$表示一次加、减、乘或除，$u$是单位舍入误差。许多局部舍入误差会在消元
中传播；带部分主元的LU通常具有良好表现，但严格后向误差界还含主元增长因子，不能
无条件声称所有矩阵上的增长都很小。QR的正交变换、正定矩阵上的Cholesky以及满足各自
前提的稳定实现，也都应在这一误差模型下讨论，而不是把精确算术恒等式直接当成浮点
保证。

\subsection{QR直接处理最小二乘}

若$\bm{A}\in\mathbb{R}^{m\times n}$且$m\geq n$、列满秩，经济型QR分解写作
$\bm{A}=\bm{Q}\bm{R}$，其中$\bm{Q}\in\mathbb{R}^{m\times n}$满足
$\bm{Q}^{\top}\bm{Q}=\bm{I}$，
$\bm{R}\in\mathbb{R}^{n\times n}$为可逆上三角矩阵。最小二乘目标变为
\[
  \lVert\bm{A}\bm{x}-\bm{b}\rVert_2^2
  =\lVert\bm{R}\bm{x}-\bm{Q}^{\top}\bm{b}\rVert_2^2+
  \lVert(\bm{I}-\bm{Q}\bm{Q}^{\top})\bm{b}\rVert_2^2.
\]
第二项与$\bm{x}$无关，所以只需求解三角系统
$\bm{R}\bm{x}=\bm{Q}^{\top}\bm{b}$。

Householder反射是计算稠密QR的常用稳定方法。经典Gram--Schmidt在列近相关时可能因
消去误差丢失正交性；矩阵可能秩亏时，可用带列主元的QR估计数值秩，或用SVD直接检查
奇异值\scite{math-golub2013matrix}

对$\bm{X}_{\widetilde{\varepsilon}}$，正规方程把条件数平方，而QR只通过正交变换和三角求解处理
原矩阵的尺度。SVD成本通常更高，却能显式识别$\sigma_{\min}$并支持截断或正则化。
选择方法时应同时考虑精度、秩判断需求与计算规模。

\subsection{Cholesky利用正定性}

若$\bm{A}$实对称正定，则存在唯一的下三角矩阵$\bm{L}$，其对角元为正，并满足
\begin{equation}
  \bm{A}=\bm{L}\bm{L}^{\top}.
  \label{eq:matrix-analysis-cholesky}
\end{equation}
这称为Cholesky分解。求解时依次计算
$\bm{L}\bm{z}=\bm{b}$和
$\bm{L}^{\top}\bm{x}=\bm{z}$。与一般LU相比，它利用对称性减少约一半存储和运算，
并且不需要主元选取。

半正定矩阵可能出现零主元，理论上半正定的矩阵也可能因舍入产生微小负特征值，导致
Cholesky失败。盲目添加固定“小常数”会改变问题；应先判断原因，再选择正则化、带主元
分解或特征值修正。

对岭回归，
$\bm{X}^{\top}\bm{X}+\lambda\bm{I}$在$\lambda>0$时正定，可以用Cholesky求解。
但形成$\bm{X}^{\top}\bm{X}$仍可能损失信息。在高精度要求或极端病态情形下，把正则化
问题写成增广最小二乘
$\min_{\bm{\beta}}\lVert
[\bm{X}^{\top},\sqrt{\lambda}\bm{I}]^{\top}\bm{\beta}
-[\bm{y}^{\top},\bm{0}^{\top}]^{\top}\rVert_2$
并用QR求解，可以避免显式平方原设计矩阵的条件数。

\subsection{稀疏大规模问题依赖矩阵向量乘法}

当$\bm{A}$规模很大且稀疏时，存储完整分解可能比原矩阵稠密得多。
\term{迭代法}（iterative method）从初值出发，通过矩阵向量乘法逐步逼近解。
对对称正定系统，共轭梯度法在精确算术中至多$d$步得到精确解；实际浮点计算中更重要的
结论是，其收敛速度受特征值分布和条件数影响。

共轭梯度的经典误差界含有因子
$((\sqrt{\kappa_2(\bm{A})}-1)/(\sqrt{\kappa_2(\bm{A})}+1))^t$：
条件数大时最坏情况收敛较慢，但特征值聚簇常使实际收敛优于只看端点的界。

\term{预条件}（preconditioning）选择容易求解的
$\bm{M}\approx\bm{A}$，把系统变换为谱更集中的等价问题。目标不是让
$\bm{M}^{-1}\bm{A}$在符号上更复杂，而是减少迭代次数，同时保持每步成本可控。
预条件器若构造或应用过于昂贵，即使降低了条件数也未必缩短总时间。非对称系统通常使用
GMRES等Krylov子空间方法；不同方法的适用条件不能互换。

\subsection{残差、前向误差与后向误差回答不同问题}

先设$\bm{A}\in\mathbb{R}^{d\times d}$可逆，
$\bm{x}=\bm{A}^{-1}\bm{b}$是原方程的唯一真解。给定计算结果
$\widehat{\bm{x}}$，其
\term{残差}（residual）为
\begin{equation}
  \bm{r}=\bm{b}-\bm{A}\widehat{\bm{x}}.
  \label{eq:matrix-analysis-residual}
\end{equation}
残差小表示$\widehat{\bm{x}}$几乎满足原方程，但不直接保证它接近真解
$\bm{x}$。因为
$\bm{x}-\widehat{\bm{x}}=\bm{A}^{-1}\bm{r}$，病态矩阵可以把很小的残差转成很大的
前向误差。

一个二阶例子足以看到差别。取
$\bm{A}=\operatorname{diag}(1,\epsilon)$、
$\bm{b}=(1,\epsilon)^{\top}$，其中$0<\epsilon\ll1$。真解为
$\bm{x}=(1,1)^{\top}$。若计算结果是
$\widehat{\bm{x}}=(1,0)^{\top}$，则
\[
  \bm{r}=(0,\epsilon)^{\top},\qquad
  \frac{\lVert\bm{r}\rVert_2}{\lVert\bm{b}\rVert_2}
  =\frac{\epsilon}{\sqrt{1+\epsilon^2}},
  \qquad
  \frac{\lVert\widehat{\bm{x}}-\bm{x}\rVert_2}
       {\lVert\bm{x}\rVert_2}
  =\frac{1}{\sqrt{2}}.
\]
相对残差随$\epsilon$趋于零，前向误差却保持常数量级，因为
$\kappa_2(\bm{A})=1/\epsilon$。小残差是方程几乎满足的证据，不是参数已经求准的
充分条件。

\term{前向误差}（forward error）直接比较
$\widehat{\bm{x}}$与真解；\term{后向误差}（backward error）则问：
$\widehat{\bm{x}}$究竟是哪一个邻近问题的精确解？若只允许扰动右端，
有$\bm{A}\widehat{\bm{x}}=\bm{b}-\bm{r}$，
所以当$\bm{b}\neq\bm{0}$时，
$\lVert\bm{r}\rVert_2/\lVert\bm{b}\rVert_2$就是只扰动右端的相对后向误差。
当$\bm{b}=\bm{0}$时，这个比值没有定义，不能把分母改成任意小常数；此时可以报告
绝对残差，或允许矩阵与右端共同扰动。

对一般矩阵$\bm{A}\in\mathbb{R}^{m\times n}$，共同扰动的范数型相对后向误差定义为
\[
  \eta_2(\widehat{\bm{x}};\bm{A},\bm{b})
  =
  \min_{\epsilon,\Delta\bm{A},\Delta\bm{b}}
  \left\{
    \epsilon\geq0:
    \begin{array}{l}
      (\bm{A}+\Delta\bm{A})\widehat{\bm{x}}
      =\bm{b}+\Delta\bm{b},\\
      \lVert\Delta\bm{A}\rVert_2\leq\epsilon\lVert\bm{A}\rVert_2,\quad
      \lVert\Delta\bm{b}\rVert_2\leq\epsilon\lVert\bm{b}\rVert_2
    \end{array}
  \right\}.
\]
令$d=\lVert\bm{A}\rVert_2\lVert\widehat{\bm{x}}\rVert_2
+\lVert\bm{b}\rVert_2$。当$d>0$时，这个最小值具有闭式
\[
  \eta_2(\widehat{\bm{x}};\bm{A},\bm{b})
  =
  \frac{\lVert\bm{r}\rVert_2}
       {\lVert\bm{A}\rVert_2\lVert\widehat{\bm{x}}\rVert_2
        +\lVert\bm{b}\rVert_2}.
\]
确实，扰动方程等价于
$\Delta\bm{A}\widehat{\bm{x}}-\Delta\bm{b}=\bm{r}$，三角不等式先给出右侧比值
作为下界；再让$\Delta\bm{A}\widehat{\bm{x}}$与
$-\Delta\bm{b}$都沿$\bm{r}$方向，并按两项分母的大小分配残差，就能达到该下界。
这个定义也覆盖$\bm{b}=\bm{0}$且$d>0$的情形，此时全部等效扰动由矩阵承担。若$d=0$，
则必有$\bm{b}=\bm{0}$且$\bm{A}\widehat{\bm{x}}=\bm{0}$，计算结果已经精确满足原
方程，约定$\eta_2=0$。

一个算法称为后向稳定，意指它返回的是某个与输入很接近的问题的精确解。后向稳定加上
问题条件良好，才能推出前向误差小。若$\kappa(\bm{A})$很大，即使算法完全后向稳定，
前向解仍可能不可靠。这一区分把责任分清：条件数描述问题对输入误差的敏感性，后向稳定性
描述算法额外引入了多大等效扰动。

若输出本身需要逆矩阵的许多元素，可以通过稳定分解求逆并检查条件数。但若目标只是计算
$\bm{A}^{-1}\bm{b}$、二次型
$\bm{b}^{\top}\bm{A}^{-1}\bm{b}$或
$\bm{A}^{-1}\bm{U}$，线性求解通常更省计算、更少存储，也更容易复用结构。对
式~\eqref{eq:matrix-analysis-nearly-collinear-design}这样的病态矩阵，换一种写法
不能恢复数据中没有的信息；合理做法是报告敏感性、引入有依据的正则化，或把目标改为
可稳定识别的预测和子空间。

\section{本章小结}
\label{sec:matrix-analysis-summary}

矩阵分析把“可逆或不可逆”的二分判断细化为方向性的描述。对称谱分解与Rayleigh商找出
二次型最强和最弱的方向；SVD把这一观点推广到任意矩阵，并同时刻画值域、零空间、低秩
近似和最小二乘。矩阵范数把不同方向的作用汇总成适合具体误差问题的尺度，有效秩则避免
把微小奇异值与主要结构等量齐观；这里的稳定秩不同于核方法中随正则化尺度变化的有效
维度。

近共线例说明了这些工具为何需要连在一起使用。小奇异值意味着某个系数方向几乎被设计
矩阵压扁；伪逆会在反向求解时放大该方向，截断SVD与岭正则化则以不同方式抑制它。由此
产生的系数敏感不必等同于训练拟合值敏感，但会削弱参数解释，并可能危及训练数据范围外
的预测。谱间隙进一步决定特征向量、奇异向量或相应子空间在扰动下是否可辨识。

可靠计算还需要把问题条件与算法稳定性分开。LU、QR、Cholesky、SVD和迭代法利用的结构
不同；残差小只说明方程满足得好，只有结合条件数与后向误差，才能判断前向解是否可信。
矩阵函数、Schur补、Woodbury恒等式、Kronecker积和极分解则展示了另一条原则：先识别
结构，再选择运算形式。矩阵函数的Fréchet导数还说明，非交换扰动不能照抄标量导数。
后续使用这些工具时，既要写出成立条件，也要说明保留、舍弃或正则化了哪些方向。
第~\ref{chap:differential-geometry-and-symmetry}章将进一步离开固定线性空间，研究
弯曲点带上的局部坐标、切空间、度量与对称性。
